\documentclass[11pt]{article}

\usepackage{amsmath,amssymb,amsfonts,amsthm}
\usepackage{mathtools}
\usepackage{bm}
\usepackage{geometry}
\geometry{margin=1in}

\title{Theoretical Analysis of One-Sided Temporal-Difference Clipping in PPO}
\author{}
\date{}

%%%%%%%%%%%%%%%%%%%%%%%%%%%%%%%%%%%%%%%%%%%%%%%%%%%%%%%%%%%%%%%

\newtheorem{assumption}{Assumption}
\newtheorem{definition}{Definition}
\newtheorem{lemma}{Lemma}
\newtheorem{theorem}{Theorem}
\newtheorem{corollary}{Corollary}
\newtheorem{remark}{Remark}

%%%%%%%%%%%%%%%%%%%%%%%%%%%%%%%%%%%%%%%%%%%%%%%%%%%%%%%%%%%%%%%

\begin{document}

\maketitle

\begin{abstract}

This appendix presents a theoretical analysis of one-sided
temporal-difference (TD) clipping applied before generalized
advantage estimation (GAE) in Proximal Policy Optimization (PPO).

Unlike conventional PPO clipping, the proposed operator modifies
the TD errors before advantage estimation.
We show that this operation introduces a sign-dependent bias into
the advantage estimator.
Under suitable assumptions, this bias propagates into the expected
policy gradient.
For binary eating decisions, sufficiently severe positive clipping
suppresses positive reinforcement whereas severe negative clipping
suppresses negative reinforcement.
Consequently, repeated retraining converges to stationary points of
a clipped optimization objective rather than those of the original
PPO objective.

\end{abstract}

%%%%%%%%%%%%%%%%%%%%%%%%%%%%%%%%%%%%%%%%%%%%%%%%%%%%%%%%%%%%%%%

\section{Preliminaries}

Consider the infinite-horizon discounted Markov decision process

\[
\mathcal{M}
=
(\mathcal S,\mathcal A,P,r,\gamma),
\]

where

\[
0<\gamma<1.
\]

The action space is binary,

\[
\mathcal A
=
\{0,1\},
\]

where

\[
a=1
\]

denotes the \emph{Eat} action and

\[
a=0
\]

denotes the \emph{Do Not Eat} action.

The stochastic policy is parameterized by

\[
\pi_\theta(a|s).
\]

Because the action space is binary,

\[
p_\theta(s)
=
\pi_\theta(1|s),
\]

and therefore

\[
\pi_\theta(0|s)
=
1-p_\theta(s).
\]

%%%%%%%%%%%%%%%%%%%%%%%%%%%%%%%%%%%%%%%%%%%%%%%%%%%%%%%%%%%%%%%

\section{Value Functions}

The state-value function is

\[
V^\pi(s)
=
\mathbb E_\pi
\left[
\sum_{k=0}^{\infty}
\gamma^k r_{t+k}
\mid
s_t=s
\right].
\]

The action-value function is

\[
Q^\pi(s,a)
=
\mathbb E_\pi
\left[
\sum_{k=0}^{\infty}
\gamma^k r_{t+k}
\mid
s_t=s,
a_t=a
\right].
\]

The advantage function is

\[
A^\pi(s,a)
=
Q^\pi(s,a)
-
V^\pi(s).
\]

%%%%%%%%%%%%%%%%%%%%%%%%%%%%%%%%%%%%%%%%%%%%%%%%%%%%%%%%%%%%%%%

\section{Temporal-Difference Errors}

The one-step temporal-difference error is

\[
\delta_t
=
r_t
+
\gamma
V(s_{t+1})
-
V(s_t).
\]

%%%%%%%%%%%%%%%%%%%%%%%%%%%%%%%%%%%%%%%%%%%%%%%%%%%%%%%%%%%%%%%

\begin{assumption}[Locally Consistent Critic]

The critic is sufficiently accurate that

\[
\mathbb E
[
\delta_t
\mid
s_t,a_t
]
=
A^\pi(s_t,a_t).
\]

\end{assumption}

This assumption is standard in analyses of actor-critic algorithms.
It states that the temporal-difference error is an unbiased estimator
of the true advantage.

%%%%%%%%%%%%%%%%%%%%%%%%%%%%%%%%%%%%%%%%%%%%%%%%%%%%%%%%%%%%%%%

\section{Generalized Advantage Estimation}

Generalized Advantage Estimation (GAE) computes

\[
A_t^{(\lambda)}
=
\sum_{l=0}^{\infty}
(\gamma\lambda)^l
\delta_{t+l},
\]

where

\[
0\le\lambda\le1.
\]

Taking expectations,

\[
\mathbb E
[
A_t^{(\lambda)}
]
=
\sum_{l=0}^{\infty}
(\gamma\lambda)^l
\mathbb E[\delta_{t+l}].
\]

Using Assumption 1,

\[
\boxed{
\mathbb E
[
A_t^{(\lambda)}
]
=
A^\pi(s_t,a_t).
}
\]

Thus standard GAE remains an unbiased estimator of the policy
advantage.

%%%%%%%%%%%%%%%%%%%%%%%%%%%%%%%%%%%%%%%%%%%%%%%%%%%%%%%%%%%%%%%

\section{One-Sided Temporal-Difference Clipping}

Instead of clipping PPO importance ratios,
consider clipping the TD errors before GAE.

\begin{definition}[One-Sided TD Clipping]

Define

\[
f_c:\mathbb R\rightarrow\mathbb R
\]

by

\[
f_c(\delta)
=
\begin{cases}
\min(\delta,c),
&
c>0,
\\[2ex]
\max(\delta,c),
&
c<0.
\end{cases}
\]

The clipped TD error is

\[
\widetilde\delta_t
=
f_c(\delta_t).
\]

The resulting clipped GAE is

\[
\widetilde A_t
=
\sum_{l=0}^{\infty}
(\gamma\lambda)^l
\widetilde\delta_{t+l}.
\]

\end{definition}

%%%%%%%%%%%%%%%%%%%%%%%%%%%%%%%%%%%%%%%%%%%%%%%%%%%%%%%%%%%%%%%

\section{Properties of One-Sided Clipping}

\begin{lemma}

Suppose

\[
c>0.
\]

Then

\[
f_c(\delta)
\le
\delta
\]

for every

\[
\delta\in\mathbb R.
\]

Furthermore,

\[
f_c(\delta)
<
\delta
\]

if and only if

\[
\delta>c.
\]

\end{lemma}

\begin{proof}

For positive clipping,

\[
f_c(\delta)
=
\min(\delta,c).
\]

By definition of the minimum,

\[
\min(\delta,c)
\le
\delta.
\]

Equality holds exactly when

\[
\delta
\le
c.
\]

Otherwise,

\[
\delta>c,
\]

which implies

\[
f_c(\delta)
=
c
<
\delta.
\]

\end{proof}

%%%%%%%%%%%%%%%%%%%%%%%%%%%%%%%%%%%%%%%%%%%%%%%%%%%%%%%%%%%%%%%

\begin{lemma}

Suppose

\[
c<0.
\]

Then

\[
f_c(\delta)
\ge
\delta
\]

for every

\[
\delta.
\]

Moreover,

\[
f_c(\delta)
>
\delta
\]

if and only if

\[
\delta<c.
\]

\end{lemma}

\begin{proof}

Negative clipping is

\[
f_c(\delta)
=
\max(\delta,c).
\]

Since the maximum of two numbers cannot be smaller than either,

\[
\max(\delta,c)
\ge
\delta.
\]

Equality holds whenever

\[
\delta
\ge
c.
\]

Otherwise,

\[
\delta<c,
\]

and therefore

\[
f_c(\delta)
=
c
>
\delta.
\]

\end{proof}

%%%%%%%%%%%%%%%%%%%%%%%%%%%%%%%%%%%%%%%%%%%%%%%%%%%%%%%%%%%%%%%

\begin{lemma}[Expectation Bias]

Suppose clipping occurs with positive probability.

For positive clipping,

\[
c>0,
\]

\[
\mathbb E[f_c(\delta)]
<
\mathbb E[\delta].
\]

For negative clipping,

\[
c<0,
\]

\[
\mathbb E[f_c(\delta)]
>
\mathbb E[\delta].
\]

\end{lemma}

\begin{proof}

From Lemma 1,

\[
f_c(\delta)
\le
\delta.
\]

Taking expectations preserves inequalities,

\[
\mathbb E[f_c(\delta)]
\le
\mathbb E[\delta].
\]

If

\[
P(\delta>c)>0,
\]

then strict inequality holds over a set of non-zero measure, implying

\[
\mathbb E[f_c(\delta)]
<
\mathbb E[\delta].
\]

The proof for negative clipping follows identically from Lemma 2.

\end{proof}

%%%%%%%%%%%%%%%%%%%%%%%%%%%%%%%%%%%%%%%%%%%%%%%%%%%%%%%%%%%%%%%

\section{Bias of Clipped GAE}

The previous lemmas immediately imply the following result.

\begin{theorem}[Bias of the Clipped Advantage Estimator]

Suppose clipping occurs with positive probability.

Then

\[
\mathbb E[\widetilde A_t]
\neq
\mathbb E[A_t].
\]

Specifically,

positive clipping satisfies

\[
\mathbb E[\widetilde A_t]
<
\mathbb E[A_t],
\]

whereas negative clipping satisfies

\[
\mathbb E[\widetilde A_t]
>
\mathbb E[A_t].
\]

\end{theorem}

\begin{proof}

Using linearity of expectation,

\[
\mathbb E[\widetilde A_t]
=
\sum_{l=0}^{\infty}
(\gamma\lambda)^l
\mathbb E[\widetilde\delta_{t+l}].
\]

Since

\[
(\gamma\lambda)^l>0,
\]

each coefficient preserves the ordering established in Lemma 3.

Therefore,

\[
\mathbb E[\widetilde A_t]
<
\mathbb E[A_t]
\]

for positive clipping and

\[
\mathbb E[\widetilde A_t]
>
\mathbb E[A_t]
\]

for negative clipping.

\end{proof}

%%%%%%%%%%%%%%%%%%%%%%%%%%%%%%%%%%%%%%%%%%%%%%%%%%%%%%%%%%%%%%%

% -------- CONTINUE IN PART II --------

%%%%%%%%%%%%%%%%%%%%%%%%%%%%%%%%%%%%%%%%%%%%%%%%%%%%%%%%%%%%%%%
\section{Policy Gradient Analysis}

The PPO objective (ignoring the probability-ratio clipping term for
clarity) is

\[
J(\theta)
=
\mathbb E_\pi
\left[
A_t
\log\pi_\theta(a_t|s_t)
\right].
\]

Its policy gradient is

\[
\nabla_\theta J(\theta)
=
\mathbb E_\pi
\left[
A_t
\nabla_\theta
\log\pi_\theta(a_t|s_t)
\right].
\]

Replacing the standard GAE estimator with the clipped estimator gives

\[
\nabla_\theta J_c(\theta)
=
\mathbb E_\pi
\left[
\widetilde A_t
\nabla_\theta
\log\pi_\theta(a_t|s_t)
\right].
\]

The remainder of this appendix analyzes the relationship between these
two gradients.

%%%%%%%%%%%%%%%%%%%%%%%%%%%%%%%%%%%%%%%%%%%%%%%%%%%%%%%%%%%%%%%

\section{Binary Bernoulli Policy}

For a binary action space,

\[
p_\theta(s)
=
\pi_\theta(a=1|s).
\]

The policy distribution is

\[
\pi_\theta(a|s)
=
p_\theta(s)^a
\left(
1-p_\theta(s)
\right)^{1-a}.
\]

Hence

\[
\log\pi_\theta(a|s)
=
a\log p_\theta(s)
+
(1-a)\log(1-p_\theta(s)).
\]

Differentiating,

\[
\nabla_\theta
\log\pi_\theta(a|s)
=
\frac{a-p_\theta(s)}
{p_\theta(s)(1-p_\theta(s))}
\nabla_\theta p_\theta(s).
\]

For the eating action,

\[
a=1,
\]

the score function becomes

\[
\nabla_\theta
\log\pi_\theta(1|s)
=
\frac{1}{p_\theta(s)}
\nabla_\theta p_\theta(s).
\]

Similarly,

\[
\nabla_\theta
\log\pi_\theta(0|s)
=
-
\frac{1}{1-p_\theta(s)}
\nabla_\theta p_\theta(s).
\]

%%%%%%%%%%%%%%%%%%%%%%%%%%%%%%%%%%%%%%%%%%%%%%%%%%%%%%%%%%%%%%%

\section{Mixed Reinforcement Assumption}

The following assumption formalizes the nutritional environment.

\begin{assumption}[Mixed Eating Outcomes]

There exists a measurable subset of eating transitions satisfying

\[
P(\delta>0\,|\,a=1)>0
\]

and

\[
P(\delta<0\,|\,a=1)>0.
\]

Thus eating may either improve or worsen the expected return,
depending upon the current physiological state.

\end{assumption}

This assumption is naturally satisfied in nutritional regulation
problems because eating may alleviate nutrient deficiency while also
causing nutrient excess.

%%%%%%%%%%%%%%%%%%%%%%%%%%%%%%%%%%%%%%%%%%%%%%%%%%%%%%%%%%%%%%%

\section{Severe Clipping}

\begin{assumption}[Severe One-Sided Clipping]

Suppose clipping is sufficiently severe that

\[
P(\delta>c)
=
1-\varepsilon,
\qquad
0<\varepsilon\ll1,
\]

for positive clipping,

or

\[
P(\delta<c)
=
1-\varepsilon,
\]

for negative clipping.

Thus almost every TD error on one side of the distribution is clipped.

\end{assumption}

%%%%%%%%%%%%%%%%%%%%%%%%%%%%%%%%%%%%%%%%%%%%%%%%%%%%%%%%%%%%%%%

\section{Gradient Propagation Assumption}

The previous theorem established that

\[
\mathbb E[\widetilde A]
\neq
\mathbb E[A].
\]

However, the policy gradient depends upon

\[
\mathbb E
[
A
\nabla\log\pi
].
\]

To relate these quantities we require one additional assumption.

\begin{assumption}[Positive Gradient Alignment]

For eating actions,

the expected score function is positively aligned with the
expected advantage, namely,

\[
\operatorname{Cov}
\left(
A_t,
\nabla_\theta
\log\pi_\theta(a_t|s_t)
\right)
\ge0.
\]

Equivalently,

decreasing the expected advantage cannot increase the expected
policy update.

\end{assumption}

This assumption is considerably weaker than independence and is
commonly satisfied near local optima in actor-critic methods.

%%%%%%%%%%%%%%%%%%%%%%%%%%%%%%%%%%%%%%%%%%%%%%%%%%%%%%%%%%%%%%%

\begin{lemma}

Under Assumption 4,

positive clipping weakens the expected policy update,

whereas negative clipping strengthens it.

\end{lemma}

\begin{proof}

Using the covariance decomposition,

\[
\mathbb E
[
A\nabla\log\pi
]
=
\mathbb E[A]
\,
\mathbb E[\nabla\log\pi]
+
\operatorname{Cov}
(A,\nabla\log\pi).
\]

Likewise,

\[
\mathbb E
[
\widetilde A\nabla\log\pi
]
=
\mathbb E[\widetilde A]
\,
\mathbb E[\nabla\log\pi]
+
\operatorname{Cov}
(
\widetilde A,
\nabla\log\pi
).
\]

By Theorem 1,

\[
\mathbb E[\widetilde A]
<
\mathbb E[A]
\]

for positive clipping.

Under Assumption 4,

the covariance term remains non-negative.

Consequently,

\[
\mathbb E
[
\widetilde A
\nabla\log\pi
]
<
\mathbb E
[
A
\nabla\log\pi
].
\]

The proof for negative clipping is identical.

\end{proof}

%%%%%%%%%%%%%%%%%%%%%%%%%%%%%%%%%%%%%%%%%%%%%%%%%%%%%%%%%%%%%%%

\begin{theorem}[Gradient Bias]

Suppose Assumptions
1--4 hold.

Then

\[
\mathbb E
[
\nabla_\theta J_c
]
\neq
\mathbb E
[
\nabla_\theta J
].
\]

Specifically,

positive clipping satisfies

\[
\mathbb E
[
\nabla_\theta J_c
]
<
\mathbb E
[
\nabla_\theta J
],
\]

whereas

negative clipping satisfies

\[
\mathbb E
[
\nabla_\theta J_c
]
>
\mathbb E
[
\nabla_\theta J
].
\]

\end{theorem}

\begin{proof}

By definition,

\[
\nabla_\theta J
=
\mathbb E
[
A
\nabla\log\pi
],
\]

while

\[
\nabla_\theta J_c
=
\mathbb E
[
\widetilde A
\nabla\log\pi
].
\]

Lemma 4 establishes

\[
\mathbb E
[
\widetilde A
\nabla\log\pi
]
<
\mathbb E
[
A
\nabla\log\pi
]
\]

under positive clipping.

The negative clipping proof follows identically.

Therefore,

the expected optimization direction differs from that of
standard PPO.

\end{proof}

%%%%%%%%%%%%%%%%%%%%%%%%%%%%%%%%%%%%%%%%%%%%%%%%%%%%%%%%%%%%%%%

\section{Retraining}

Repeated PPO updates satisfy

\[
\theta_{k+1}
=
\theta_k
+
\alpha
\nabla_\theta J_c(\theta_k),
\]

where

\[
\alpha>0
\]

is the learning rate.

Since every update is computed from the clipped advantage estimator,

the optimization repeatedly follows the biased gradient established
above.

%%%%%%%%%%%%%%%%%%%%%%%%%%%%%%%%%%%%%%%%%%%%%%%%%%%%%%%%%%%%%%%

\begin{theorem}[Persistence Under Retraining]

Suppose clipping remains active throughout optimization.

Then repeated retraining converges, if convergence occurs,
to stationary points satisfying

\[
\nabla_\theta J_c(\theta)=0,
\]

rather than

\[
\nabla_\theta J(\theta)=0.
\]

\end{theorem}

\begin{proof}

The optimization algorithm never evaluates

\[
\nabla_\theta J.
\]

Instead,

every iteration computes

\[
\nabla_\theta J_c.
\]

Therefore gradient ascent optimizes the clipped objective

\[
J_c,
\]

not the original PPO objective.

Consequently,

all accumulation points satisfy

\[
\nabla_\theta J_c(\theta)=0.
\]

No update step can remove the clipping bias while clipping remains
active.

\end{proof}

%%%%%%%%%%%%%%%%%%%%%%%%%%%%%%%%%%%%%%%%%%%%%%%%%%%%%%%%%%%%%%%

%----------- CONTINUE IN PART III ------------------

%%%%%%%%%%%%%%%%%%%%%%%%%%%%%%%%%%%%%%%%%%%%%%%%%%%%%%%%%%%%%%%
\section{Asymptotic Behaviour of Binary Eating Policies}

The previous sections established that one-sided TD clipping
introduces a systematic bias into the expected PPO policy gradient.

We now investigate how this gradient bias affects the probability
of selecting the eating action.

%%%%%%%%%%%%%%%%%%%%%%%%%%%%%%%%%%%%%%%%%%%%%%%%%%%%%%%%%%%%%%%

\begin{assumption}[Positive Policy Responsiveness]

Suppose the Bernoulli policy satisfies

\[
\left\|
\frac{\partial p_\theta(s)}
{\partial\theta}
\right\|
>0
\]

for every state encountered during training.

Consequently,

a non-zero policy gradient induces a non-zero change in the
probability of selecting the eating action.

\end{assumption}

%%%%%%%%%%%%%%%%%%%%%%%%%%%%%%%%%%%%%%%%%%%%%%%%%%%%%%%%%%%%%%%

\begin{assumption}[Gradient Descent Stability]

Suppose the learning rate satisfies

\[
0<\alpha<\alpha^\star,
\]

where

\[
\alpha^\star
\]

is sufficiently small that stochastic gradient ascent converges
towards a stationary point of the optimization objective.

\end{assumption}

%%%%%%%%%%%%%%%%%%%%%%%%%%%%%%%%%%%%%%%%%%%%%%%%%%%%%%%%%%%%%%%

\begin{lemma}[Monotonicity of Binary Policy Updates]

Suppose

\[
a=1
\]

corresponds to eating.

Then the Bernoulli parameter satisfies

\[
p_{\theta_{k+1}}(s)
=
p_{\theta_k}(s)
+
\alpha
\frac{\partial p_\theta(s)}
{\partial\theta}
\nabla_\theta J_c
+
O(\alpha^2).
\]

\end{lemma}

\begin{proof}

Using a first-order Taylor expansion,

\[
p_{\theta_{k+1}}
=
p_{\theta_k}
+
\nabla_\theta p_\theta
\cdot
(\theta_{k+1}-\theta_k)
+
O(\alpha^2).
\]

Since PPO performs

\[
\theta_{k+1}
=
\theta_k
+
\alpha
\nabla_\theta J_c,
\]

substitution immediately yields

\[
p_{\theta_{k+1}}
=
p_{\theta_k}
+
\alpha
\nabla_\theta p_\theta
\cdot
\nabla_\theta J_c
+
O(\alpha^2).
\]

\end{proof}

%%%%%%%%%%%%%%%%%%%%%%%%%%%%%%%%%%%%%%%%%%%%%%%%%%%%%%%%%%%%%%%

\begin{theorem}[Positive One-Sided Clipping]

Suppose

\begin{enumerate}

\item Assumptions 1--6 hold,

\item eating transitions generate both positive and negative
temporal-difference errors,

\item positive clipping is sufficiently severe that

\[
P(\delta>c)
=
1-\varepsilon,
\qquad
\varepsilon\ll1.
\]

\end{enumerate}

Then

\[
\mathbb E
[
p_{\theta_{k+1}}(s)
]
<
\mathbb E
[
p_{\theta_k}(s)
]
\]

for every gradient iteration.

Consequently,

the sequence

\[
\{
p_{\theta_k}
\}
\]

is monotonically decreasing in expectation.

\end{theorem}

\begin{proof}

Theorem 2 established that

\[
\mathbb E
[
\nabla_\theta J_c
]
<
\mathbb E
[
\nabla_\theta J
].
\]

Because eating receives both positive and negative TD errors,
positive clipping suppresses only the positive reinforcement.

Negative reinforcement remains unchanged.

Therefore,

the expected policy update toward eating is strictly weaker
than that produced by standard PPO.

Applying Lemma 5,

\[
\mathbb E
[
p_{\theta_{k+1}}
]
=
\mathbb E
[
p_{\theta_k}
]
+
\alpha
\mathbb E
[
\nabla_\theta p_\theta
\cdot
\nabla_\theta J_c
]
+
O(\alpha^2).
\]

Since

\[
\nabla_\theta p_\theta
\]

is non-zero by Assumption 5,

and

\[
\nabla_\theta J_c
<
\nabla_\theta J,
\]

the first-order increment is negative.

Hence

\[
\boxed{
\mathbb E
[
p_{\theta_{k+1}}
]
<
\mathbb E
[
p_{\theta_k}
].
}
\]

Thus the probability of selecting the eating action decreases
monotonically during optimization.

\end{proof}

%%%%%%%%%%%%%%%%%%%%%%%%%%%%%%%%%%%%%%%%%%%%%%%%%%%%%%%%%%%%%%%

\begin{theorem}[Negative One-Sided Clipping]

Suppose

negative clipping satisfies

\[
P(\delta<c)
=
1-\varepsilon.
\]

Then

\[
\mathbb E
[
p_{\theta_{k+1}}
]
>
\mathbb E
[
p_{\theta_k}
].
\]

Consequently,

the eating probability increases monotonically throughout
optimization.

\end{theorem}

\begin{proof}

Negative clipping removes only large negative TD errors.

Positive TD errors remain unchanged.

Therefore

\[
\mathbb E
[
\widetilde A
]
>
\mathbb E
[
A
].
\]

By Theorem 2,

the expected policy gradient becomes larger than that of
standard PPO.

Applying Lemma 5,

the first-order Taylor expansion yields

\[
\mathbb E
[
p_{\theta_{k+1}}
]
>
\mathbb E
[
p_{\theta_k}
].
\]

Hence the Bernoulli parameter increases monotonically.

\end{proof}

%%%%%%%%%%%%%%%%%%%%%%%%%%%%%%%%%%%%%%%%%%%%%%%%%%%%%%%%%%%%%%%

\begin{corollary}[Persistent Under-Eating]

Under severe positive clipping,

repeated PPO retraining converges to a stationary policy satisfying

\[
p^\star_c
<
p^\star,
\]

where

\[
p^\star
\]

denotes the stationary eating probability obtained without
TD clipping.

\end{corollary}

\begin{proof}

Theorem 3 establishes that optimization converges to a stationary
point of

\[
J_c.
\]

Theorem 4 establishes that every optimization step decreases
the Bernoulli parameter.

Therefore,

the limiting stationary point satisfies

\[
p^\star_c
<
p^\star.
\]

\end{proof}

%%%%%%%%%%%%%%%%%%%%%%%%%%%%%%%%%%%%%%%%%%%%%%%%%%%%%%%%%%%%%%%

\begin{corollary}[Persistent Over-Eating]

Under severe negative clipping,

the stationary Bernoulli parameter satisfies

\[
p^\star_c
>
p^\star.
\]

\end{corollary}

\begin{proof}

This follows directly from Theorems 3 and 5.

\end{proof}

%%%%%%%%%%%%%%%%%%%%%%%%%%%%%%%%%%%%%%%%%%%%%%%%%%%%%%%%%%%%%%%

\section{Discussion}

The preceding analysis demonstrates that one-sided temporal-difference
clipping fundamentally alters the optimization objective used by PPO.

Unlike conventional PPO clipping, which operates on importance ratios,
the proposed operator modifies the temporal-difference signal before
advantage estimation.

Consequently,

the resulting gradient estimator is biased.

Under the mixed-reward assumption appropriate for nutritional
regulation tasks, positive clipping preferentially suppresses
positive reinforcement, while negative clipping preferentially
suppresses negative reinforcement.

Repeated retraining therefore converges to stationary policies that
systematically under-select or over-select the eating action,
respectively.

\hfill$\square$

\end{document}